\chapter{Prakata}

Buku ini ditulis untuk membantu pembaca yang belum pernah memprogram agar mampu menulis program Python sendiri, dengan AI generatif atau GenAI sebagai pendamping belajar yang terarah. Pembahasannya dimulai dari program satu baris, lalu berlanjut sampai aplikasi berantarmuka grafis yang diuji dan diukur kinerjanya. Buku ini menjadi pegangan mata kuliah Pemrograman Dasar dan Praktikum Pemrograman Dasar di Program Studi Informatika, Universitas Pradita.

Layanan GenAI seperti ChatGPT, Gemini, dan Claude dapat menjelaskan kode, menebak keluaran program, bahkan menuliskan program utuh dalam beberapa detik. Kemudahan ini dapat merugikan pemula bila dipakai tanpa arah. Dalam eksperimen lapangan dengan hampir seribu siswa SMA, siswa yang memakai GPT-4 tanpa pengaman mendapat nilai latihan lebih tinggi, tetapi setelah aksesnya dicabut nilai ujian mereka 17\% lebih rendah daripada siswa yang tidak pernah memakainya. Penurunan itu sebagian besar tidak terjadi pada siswa yang memakai GPT-4 dengan \textit{prompt} pengaman \cite{bastani2025}. Buku ini menyebut cara pakai seperti itu GenAI terarah. Pembaca memberi GenAI peran dan batasan lewat \textit{prompt}, misalnya menjelaskan kode tanpa menulis ulang, memberi petunjuk tanpa jawaban jadi, atau memeriksa jawaban pembaca. Kode tetap ditulis dan dijalankan sendiri, dan penjelasan GenAI dicocokkan dengan hasil menjalankan kode karena GenAI bisa keliru.

Pembaca sasarannya adalah mahasiswa tahun pertama program studi informatika dan bidang yang berdekatan. Pembaca umum yang belajar mandiri juga dapat memakainya, karena tidak ada pengalaman memprogram yang diandaikan. Yang dibutuhkan hanya komputer yang dapat dipasangi Python. Contoh kode ditulis untuk Python 3.10 atau lebih baru.

Isi buku terbagi menjadi enam belas bab. Bab 1 mengenalkan Python dan cara menyiapkan lingkungan kerja. Bab 2 sampai Bab 5 membahas dasar pemrograman, yaitu variabel, tipe data, operator, percabangan, fungsi, modul, dan perulangan. Bab 6 dan Bab 7 membahas cara mengolah kumpulan data dan berkas. Bab 8 mengenalkan pemrograman berorientasi objek. Bab 9, Bab 10, dan Bab 13 membahas cara menangani, menguji, dan menelusuri kesalahan. Bab 11 dan Bab 12 membahas antarmuka grafis dengan Tkinter, Bab 14 membahas pengukuran kinerja, dan Bab 15 berisi \textit{capstone project} membuat \textit{game} yang memakai seluruh materi. Bab 16 melaporkan penelitian tentang hubungan cara mahasiswa memakai GenAI dengan keyakinan mereka bisa memprogram. Bab itu juga memuat dua \textit{prompt} umum untuk meminta penjelasan kode dan soal latihan baru.

Tiap bab materi dibuka dengan tujuan pembelajaran dan petunjuk memakai GenAI terarah untuk topik bab itu, lalu ditutup dengan ringkasan, latihan, dan referensi. Setiap soal latihan disertai \textit{prompt} terarah yang meminta GenAI memberi petunjuk tanpa jawaban jadi. Banyak contoh kode disertai keluaran dari hasil menjalankannya, sehingga pembaca dapat mencocokkan hasilnya sendiri. Kesalahan yang sering dibuat pemula ikut dibahas beserta pesan kesalahan yang muncul. Rujukan diutamakan dari dokumentasi resmi dan buku terbuka yang dapat dibuka secara daring. Materi buku ini juga berasal dari modul yang sudah dipakai di kelas.

Cara terbaik memakai buku ini adalah mengetik dan menjalankan sendiri setiap contoh, lalu mengubahnya untuk melihat apa yang terjadi. Bila sebuah listing belum dipahami, salin listing itu ke layanan GenAI bersama \textit{prompt} pada Listing~\ref{lst:b16-prompt-penjelasan}, lalu jawab tiga pertanyaan yang diajukannya. Pesan kesalahan bukan tanda kegagalan, karena pesan itu menunjukkan di mana program perlu diperbaiki. Kode latihan untuk sebagian bab dan \textit{slide} pendukungnya tersedia di GitHub pada alamat \url{https://github.com/alfa-yohannis/python-programming-fundamentals}. Saran dan koreksi dari pembaca dapat disampaikan melalui penerbit.

\bigskip
\begin{flushright}
Tangerang, \isian{bulan dan tahun terbit}\\[4pt]
Penulis
\end{flushright}
